% © 2025 Ing. David Jaroš — CC BY-NC-ND 4.0
%
% This work is licensed under a Creative Commons Attribution-NonCommercial-NoDerivatives
% 4.0 International License (CC BY-NC-ND 4.0).
%
% License History: Earlier drafts (up to v0.3) were released under CC BY 4.0.
% From v0.4 onward, all material is released under CC BY-NC-ND 4.0 to protect
% the integrity of the theoretical work during ongoing academic development.

\documentclass[12pt]{article}
\usepackage{amsmath,amssymb,amsthm}
\usepackage{geometry}
\usepackage{hyperref}
\usepackage{tcolorbox}
\usepackage{xcolor}
\geometry{margin=1in}

% Theorem-like environments
\newtheorem{conjecture}{Conjecture}
\newtheorem{proposition}{Proposition}
\newtheorem{definition}{Definition}
\newtheorem{remark}{Remark}

\title{ST-2: Rotation in the $(t,\psi)$ Plane and Symmetries of\\
the UBT Field Equations\\[0.5em]
\large\textsc{Speculative Extension}}
\author{David Jaroš}
\date{2026-03-01}

\begin{document}
\maketitle

\begin{tcolorbox}[colback=yellow!10, title=Status]
\textbf{SPECULATIVE EXTENSION} — This content is not part of
the canonical UBT framework. It is an exploratory investigation
and has not been validated against the canonical field equations.
\end{tcolorbox}

\vspace{1em}
\noindent\textit{Output of Sub-task ST-2 from the CTC/Causal-Structure investigation track.
See also: \texttt{complex\_time\_causal\_structure.tex},
\texttt{ctc\_conditions.tex},
\texttt{chronology\_protection\_ubt.tex}, \texttt{ctc\_methodology.tex}.}

%----------------------------------------------------------------------
\section*{Abstract}
%----------------------------------------------------------------------
We investigate the one-parameter family of $SO(2)$ rotations in the
complex-time plane $(t,\psi)$ acting on the UBT master field $\Theta$.
The three special cases $\alpha = \pi/2$ (Wick rotation), general
$\alpha$ (complex-time rotation), and $\alpha = \pi$ (full inversion)
are examined.  We ask whether any of these are symmetries of the
canonical UBT field equations, gauge transformations, or physically
observable operations.  All conclusions are \textbf{speculative}.

%----------------------------------------------------------------------
\section{The $(t,\psi)$ Rotation}
%----------------------------------------------------------------------
Define the one-parameter rotation acting on the real and imaginary parts
of complex time $\tau = t + i\psi$:
\begin{equation}
  R_\alpha:\quad
  \begin{pmatrix} t \\ \psi \end{pmatrix}
  \;\longmapsto\;
  \begin{pmatrix} t\cos\alpha - \psi\sin\alpha \\
                  t\sin\alpha + \psi\cos\alpha \end{pmatrix},
  \qquad \alpha \in [0, 2\pi).
  \label{eq:rotation}
\end{equation}
Equivalently, on $\tau \in \mathbb{C}$:
\begin{equation}
  R_\alpha(\tau) \;=\; e^{i\alpha}\,\tau.
  \label{eq:complex_rotation}
\end{equation}
The field is transformed as
\begin{equation}
  \Theta(t,\psi,x^k) \;\longmapsto\;
  \Theta\!\left(t\cos\alpha - \psi\sin\alpha,\;
                t\sin\alpha + \psi\cos\alpha,\; x^k\right).
  \label{eq:field_transform}
\end{equation}

%----------------------------------------------------------------------
\section{Case $\alpha = \pi/2$: Wick Rotation}
%----------------------------------------------------------------------
For $\alpha = \pi/2$:
\begin{equation}
  R_{\pi/2}:\quad (t,\psi) \;\mapsto\; (-\psi,\, t).
\end{equation}
Under this map real time $t$ and imaginary time $\psi$ are exchanged
(up to a sign): $t' = -\psi$, $\psi' = t$.  On the complex plane this
is multiplication by $i$: $\tau \mapsto i\tau$.

\paragraph{Is this a symmetry of the UBT field equations?}
The canonical UBT field equation is
\begin{equation}
  \mathcal{E}_{\mu\nu}[\Theta, \mathcal{G}]
  \;=\; \kappa\,\mathcal{T}_{\mu\nu}[\Theta, \mathcal{G}].
  \label{eq:field_eq}
\end{equation}
The real projection of the metric $g_{\mu\nu} = \operatorname{Re}(\mathcal{G}_{\mu\nu})$
couples to ordinary matter.  Under Wick rotation, the Lorentzian signature
$(-,+,+,+)$ of $g_{\mu\nu}$ maps to Euclidean signature $(+,+,+,+)$.
This is not an invariance of the field equations in Lorentzian signature;
it is instead an \emph{analytic continuation} used as a computational tool.

\begin{conjecture}[Wick rotation as analytic continuation]
\label{conj:wick}
The map $R_{\pi/2}$ is not a physical symmetry of the UBT field equations
in Lorentzian signature.  It is an \emph{analytic continuation} that maps
UBT solutions in Lorentzian signature to corresponding solutions in
Euclidean signature, provided the field $\Theta$ is analytic in $\tau$ in
a strip containing the real axis.
\end{conjecture}

\begin{remark}
In standard quantum field theory the Wick rotation is used to define the
Euclidean path integral and is not a physical symmetry.  The same
interpretation is expected to hold in UBT.  The existence of the imaginary
component $\psi$ in UBT does not alter this: $\psi$ plays the role of the
Euclidean time after a Wick rotation of the real time $t$.
\end{remark}

%----------------------------------------------------------------------
\section{General $\alpha$: Physical Interpretation}
%----------------------------------------------------------------------
For a general angle $\alpha$, the rotation $R_\alpha$ mixes $t$ and $\psi$
continuously.  We consider three possibilities:

\subsection{$R_\alpha$ as a Physical Symmetry}
If the action functional $S[\Theta]$ is invariant under $R_\alpha$ for all
$\alpha$, then by Noether's theorem there is a conserved charge associated
with rotations in the $(t,\psi)$ plane.

The UBT Lagrangian density contains terms of the form
$\partial_\tau \Theta \cdot \partial_{\bar\tau} \Theta^*$ which are
invariant under $\tau \mapsto e^{i\alpha}\tau$ only if $\Theta$ transforms
with a compensating phase.  A minimal coupling
\begin{equation}
  \Theta \;\longmapsto\; e^{-i n\alpha}\,\Theta, \qquad n \in \mathbb{Z},
\end{equation}
would restore invariance.  This would constitute a $U(1)$ symmetry in
complex time and would generate a conserved ``temporal charge''.

\begin{conjecture}[$U(1)_\tau$ symmetry]
\label{conj:u1_tau}
The UBT field equations admit a $U(1)$ symmetry
$(\tau, \Theta) \mapsto (e^{i\alpha}\tau, e^{-i\alpha}\Theta)$
if and only if the biquaternionic mass term and interaction vertices are
invariant under simultaneous rescaling of $\tau$ and $\Theta$.  This
symmetry, if it exists, generates a conserved temporal charge $Q_\tau$.
Verification requires explicit computation of the UBT action functional.
\end{conjecture}

\subsection{$R_\alpha$ as a Gauge Transformation}
Alternatively, $R_\alpha$ may be a \emph{gauge redundancy}: two field
configurations related by $R_\alpha$ describe the same physical state.
In this case the phase $\alpha$ labels a gauge orbit and the physical
Hilbert space is the quotient of the full configuration space by
$R_\alpha$.

\begin{remark}
If $R_\alpha$ is a gauge transformation, then the imaginary component
$\psi$ is pure gauge and carries no independent physical information.
This aligns with Interpretation C (auxiliary degree of freedom) in
\texttt{complex\_time\_causal\_structure.tex}.
\end{remark}

\subsection{$R_\alpha$ as a Non-Symmetry Transformation}
If the UBT action changes non-trivially under $R_\alpha$, then the
rotation is neither a symmetry nor a gauge redundancy.  Different values
of $\alpha$ label genuinely distinct physical configurations.

%----------------------------------------------------------------------
\section{Case $\alpha = \pi$: Full Inversion and CPT}
%----------------------------------------------------------------------
For $\alpha = \pi$:
\begin{equation}
  R_\pi:\quad \tau \;\mapsto\; -\tau, \quad
  \text{i.e.,}\quad t \mapsto -t,\; \psi \mapsto -\psi.
\end{equation}
This simultaneously reverses real time and negates the phase coordinate.

\paragraph{Relationship to CPT.}
In standard QFT, CPT symmetry combines charge conjugation ($C$), parity
($P$), and time reversal ($T$).  The $T$-operation sends $t \mapsto -t$.
In UBT the natural extension of $T$ is:
\begin{equation}
  T_{\mathrm{UBT}}:\quad (t,\psi) \;\mapsto\; (-t, -\psi),
\end{equation}
which is precisely $R_\pi$.  If charge conjugation $C$ and parity $P$
act in the standard way on the spatial and internal degrees of freedom,
then $CPT_{\mathrm{UBT}} = C \cdot P \cdot R_\pi$.

\begin{conjecture}[CPT in UBT]
\label{conj:cpt}
The operation $R_\pi$ on complex time $\tau$ is the UBT analogue of
$T$-reversal.  Combined with $C$ and $P$ in the standard way, $R_\pi$
extends CPT symmetry to the complex-time sector.  In particular, if
the UBT Lagrangian is CPT-invariant in the real sector, it is also
invariant under $R_\pi$ in the complex sector, provided the imaginary
metric components $h_{\mu\nu}$, $a_{\mu\nu}$, $b_{\mu\nu}$ transform
covariantly.
\end{conjecture}

\begin{remark}
An important subtlety: $T$-reversal in the real sector maps Lorentzian
solutions to Lorentzian solutions.  The $\psi \mapsto -\psi$ part of
$R_\pi$ maps phase-advancing solutions to phase-receding solutions.
Whether these are physically equivalent depends on the topology of the
$\psi$ coordinate (compact or non-compact).
\end{remark}

%----------------------------------------------------------------------
\section{Summary Table}
%----------------------------------------------------------------------
\begin{center}
\renewcommand{\arraystretch}{1.4}
\begin{tabular}{|l|l|l|l|}
\hline
\textbf{Angle $\alpha$} & \textbf{Map} & \textbf{Interpretation} & \textbf{Status}\\
\hline
$\pi/2$ & $\tau \mapsto i\tau$ & Wick rotation & Analytic continuation\\
General & $\tau \mapsto e^{i\alpha}\tau$ & $U(1)_\tau$ rotation & Possible symmetry (conjecture)\\
$\pi$ & $\tau \mapsto -\tau$ & Full inversion & Candidate $T$-reversal in UBT\\
\hline
\end{tabular}
\end{center}

%----------------------------------------------------------------------
\section{Conclusions}
%----------------------------------------------------------------------
\begin{itemize}
  \item The Wick rotation ($\alpha=\pi/2$) is an analytic continuation
    tool, not a physical symmetry of UBT in Lorentzian signature.
  \item General rotations in $(t,\psi)$ may constitute a $U(1)_\tau$
    symmetry if the UBT action is invariant under simultaneous phase
    rotations of $\tau$ and $\Theta$.  This remains unverified.
  \item The full inversion $R_\pi$ is the natural extension of
    $T$-reversal to UBT's complex time, and the combined $CPT_{\mathrm{UBT}}$
    symmetry is expected to hold under standard assumptions.
\end{itemize}

\medskip
\noindent\textbf{All statements in this document are speculative.  No
claim constitutes an established result of the canonical UBT framework.}

%----------------------------------------------------------------------
\bibliographystyle{plain}
\begin{thebibliography}{9}
\bibitem{Streater1964}
  R.~F.~Streater and A.~S.~Wightman,
  \emph{PCT, Spin and Statistics, and All That},
  Princeton University Press (1964).
\bibitem{UBT_core}
  D.~Jaroš,
  \emph{Unified Biquaternion Theory: core field equations},
  Repository document \texttt{consolidation\_project/ubt\_core\_main.tex} (2025).
\bibitem{Causal_ST1}
  D.~Jaroš,
  \emph{ST-1: Causal Structure of Complex Time},
  Repository document
  \texttt{speculative\_extensions/causality/complex\_time\_causal\_structure.tex} (2026).
\end{thebibliography}

\end{document}
